\chapter{معماری شناختی، حافظه و بار شناختی}

\section{معماری شناختی}

\subsection{چرا معماری شناختی لازم است؟}
پیش از ورود به جزئیات معماری‌های شناختی، یک مثال ساده نشان می‌دهد چرا صرفاً داشتن حسگر، قانون و خروجی درست کافی نیست. اگر بخواهیم مدل هوش مصنوعی به مهارتی نزدیک شود که انسان در زندگی روزمره به‌سادگی انجام می‌دهد، باید بدانیم چه اجزایی در فرایند شناختی دخیل‌اند: ادراک، حافظه، دانش عامه، پیش‌بینی پیامد، تصمیم‌گیری، عمل و دریافت بازخورد.

معماری شناختی دقیقاً برای همین لازم می‌شود. معماری شناختی فقط فهرستی از قابلیت‌ها نیست؛ بلکه باید روشن کند اطلاعات از کجا وارد می‌شود، در کجا نگه داشته می‌شود، با کدام دانش قبلی ترکیب می‌شود، چگونه تصمیم ساخته می‌شود و نتیجهٔ عمل چگونه به سیستم برمی‌گردد.

\subsection{مثال جارو رباتیک}
مثال مناسب، جاروبرقی رباتیک است. یک ربات ساده ممکن است چند حسگر و چند قانون داشته باشد:
\begin{itemize}
  \item اگر به گردوغبار رسید، آن را جمع کند.
  \item اگر مانعی دید، مسیرش را تغییر دهد.
  \item اگر به یک شیء کوچک رسید، آن را کنار بزند.
\end{itemize}

این قواعد در بسیاری از موقعیت‌های ساده کافی‌اند. اما اگر ربات به فنجانی پر از مایع برسد، مسئله تغییر می‌کند. انسان معمولاً فنجان را بدون فکر هل نمی‌دهد، چون چند رابطهٔ علّی و پیامدی را می‌فهمد: فنجان ممکن است پر باشد؛ اگر هل داده شود، مایع می‌ریزد؛ اگر مایع بریزد، زمین خیس می‌شود؛ خیس شدن زمین نامطلوب است؛ پس نباید فنجان را مانند یک مانع معمولی کنار زد.

این مثال نشان می‌دهد مسئله فقط تشخیص شیء نیست. ربات باید چیزی شبیه دانش عامه، پیش‌بینی پیامد، ارزیابی ریسک و انتخاب عمل داشته باشد. بدون این مهارت‌ها، سیستم ممکن است در معیارهای سادهٔ مهندسی خوب به نظر برسد، اما در موقعیت‌های انسانی و پرجزئیات رفتار نامناسبی تولید کند.

\subsubsection{راه‌حل اول: دوری از اشیای پرریسک}
یک راه ساده این است که ربات اشیای پرریسک را نادیده بگیرد یا از آن‌ها دور شود. این راه‌حل خسارت را کم می‌کند، اما مسئله را کامل حل نمی‌کند. اگر ربات هر شیء مشکوکی را رها کند، ممکن است بخش‌هایی از محیط هرگز تمیز نشوند. بنابراین با یک \eng{trade-off} روبه‌رو هستیم: کاهش خطر در برابر ناقص‌ماندن وظیفه.

\subsubsection{راه‌حل دوم: تعریف تابع ریسک}
راه دیگر این است که برای حرکت دادن هر شیء، یک تابع ریسک تعریف کنیم. مدل می‌تواند بپرسد: شیء چیست؟ ظرف است یا نه؟ شکستنی است یا نه؟ داخل آن مایع هست یا نه؟ احتمال ریختن چقدر است؟ پیامد ریختن چیست؟

این مسیر از نظر مهندسی ممکن است مفید باشد، اما اگر فقط به افزودن دائمی \eng{if-else}های تازه منجر شود، سیستم شکننده و پرهزینه می‌شود. هر موقعیت تازه ممکن است قانون تازه‌ای بخواهد.

\subsubsection{راه‌حل سوم: یادگیری از تجربه و بازخورد}
راه سوم این است که ربات از تجربه و \eng{feedback} یاد بگیرد. کودک نیز از ابتدا همهٔ قوانین فیزیکی و اجتماعی را به‌صورت صریح ندارد؛ تجربه می‌کند، پیامدها را می‌بیند و از بازخورد مثبت یا منفی یاد می‌گیرد. اگر فنجان بیفتد و آب بریزد، این تجربه می‌تواند در رفتارهای بعدی اثر بگذارد.

از نگاه علوم شناختی محاسباتی، پرسش مهم این است که این یادگیری فقط حفظ نمونه‌ها نباشد، بلکه به فهم روابط علّی و تعمیم به موقعیت‌های تازه نزدیک شود.

\subsubsection{دو مسیر برای بهبود مدل}
از این مثال دو مسیر کلی برای بهبود مدل‌های هوش مصنوعی به دست می‌آید:
\begin{enumerate}
  \item مدل را بزرگ‌تر کنیم و قوانین یا داده‌های بیشتری به آن بدهیم.
  \item یک مهارت یا ماژول شناختی به مدل اضافه کنیم.
\end{enumerate}

مسیر اول می‌تواند در عمل مفید باشد، اما همیشه خطر دارد که فقط به انباشتن قوانین سطحی تبدیل شود. مسیر دوم می‌پرسد کدام توانایی شناختی کم است: فهم علّیت، حافظه، توجه، پیش‌بینی پیامد، ارزیابی ریسک یا یادگیری از بازخورد؟ این همان نقطه‌ای است که علوم شناختی به طراحی مدل هوش مصنوعی وصل می‌شود.

\subsection{ماژول‌های اصلی در یک معماری شناختی}
از مثال جارو رباتیک می‌توان اجزای اولیهٔ یک معماری شناختی ساده را استخراج کرد. ربات فقط یک ورودی و خروجی ندارد؛ میان ورودی و خروجی چند سطح پردازش لازم است. ابتدا باید محیط دیده یا حس شود، سپس وضعیت فعلی در حافظهٔ موقت قرار بگیرد، بعد دانش مرتبط از حافظهٔ بلندمدت بازیابی شود، سپس یک قاعده یا مهارت مناسب انتخاب شود و در نهایت دستور حرکتی تولید گردد.

معماری شناختی یعنی همین سازمان‌دهی یکپارچهٔ فرایندها. اگر چند نظریهٔ شناختی جداگانه داشته باشیم، معماری مشخص می‌کند این نظریه‌ها چگونه در یک سیستم واحد با هم تعامل می‌کنند: اطلاعات از کجا وارد می‌شود، کدام ماژول آن را پردازش می‌کند، چه چیزی در حافظه می‌ماند، چه چیزی بازیابی می‌شود و عمل چگونه تولید می‌شود.

\subsubsection{ادراک و ماژول بینایی}
ماژول بینایی یا ادراکی، ورودی محیط را به بازنمایی قابل استفاده برای سیستم تبدیل می‌کند. در مثال ربات، این ماژول باید تشخیص دهد که یک فنجان وجود دارد، داخل آن احتمالاً آب است و فنجان در مسیر حرکت قرار گرفته است.

در نگاشت تقریبی به مغز، پردازش بینایی با نواحی بینایی، به‌ویژه لوب پس‌سری، ارتباط دارد. البته در مغز واقعی، بینایی فقط در یک نقطهٔ جداگانه انجام نمی‌شود و مسیرهای گسترده‌تری برای تشخیص شیء، مکان، حرکت و معنا درگیرند. اما برای سطح فعلی معماری، می‌توان ماژول بینایی را ورودی ادراکی سیستم دانست.

\subsubsection{حافظه و بافرهای کاری}
اطلاعاتی که از محیط می‌آید باید موقتاً نگه داشته شود. در معماری‌های شناختی معمولاً از مفهوم \term{بافر}{Buffer} استفاده می‌شود: فضای محدودی که وضعیت فعلی سیستم را در دسترس سایر ماژول‌ها قرار می‌دهد. بافر را می‌توان به حافظهٔ کوتاه‌مدت یا حافظهٔ کاری نزدیک دانست، هرچند در معماری‌هایی مانند \eng{ACT-R} معنای فنی دقیق‌تری دارد.

برای مثال، پس از دیدن فنجان، محتوای بافر می‌تواند چیزی شبیه این باشد: «فنجان در مسیر است» و «مایع داخل فنجان دیده می‌شود». این اطلاعات هنوز به‌تنهایی تصمیم ایجاد نمی‌کند؛ فقط وضعیت فعلی محیط را برای بازیابی دانش و انتخاب عمل آماده می‌کند.

\subsubsection{تصمیم‌گیری و انتخاب عمل}
پس از قرار گرفتن وضعیت فعلی در بافر، سیستم باید تصمیم بگیرد چه عملی مناسب است. این تصمیم از ترکیب چند منبع به دست می‌آید: وضعیت فعلی، دانش اعلانی دربارهٔ اشیا و پیامدها، و قواعد یا مهارت‌های رویه‌ای.

اگر سیستم بداند فنجان شکستنی است، آب ممکن است بریزد و خیس شدن زمین نامطلوب است، آنگاه یک قاعدهٔ رویه‌ای می‌تواند فعال شود: اگر فنجان پر از آب در مسیر است، با آن مثل مانع معمولی برخورد نکن. این مرحله همان جایی است که تفاوت یک مدل صرفاً واکنشی با یک مدل شناختی‌تر آشکار می‌شود.

\subsubsection{ماژول حرکتی}
ماژول حرکتی تصمیم مفهومی را به عمل فیزیکی تبدیل می‌کند. «تغییر مسیر بده» باید به دستورهایی مانند کاهش سرعت، چرخش، حرکت بازو، گرفتن آرام فنجان یا بازگشت به مسیر قبلی تبدیل شود.

در نگاشت مغزی، کنترل حرکت با بخش‌هایی از لوب پیشانی و مدارهای حرکتی مرتبط است. باید میان \eng{frontal lobe} به‌عنوان ناحیهٔ وسیع‌تر و \eng{prefrontal cortex} به‌عنوان بخشی که بیشتر با برنامه‌ریزی، استدلال، کنترل اجرایی و تصمیم‌گیری مرتبط است تمایز گذاشت.

\subsection{نگاشت معماری شناختی به مغز}
یکی از شروط مهم معماری شناختی این است که اجزای آن دست‌کم به‌صورت تقریبی با مغز قابل نگاشت باشند. اگر معماری یک ماژول بینایی، ماژول حرکتی، حافظهٔ اعلانی و حافظهٔ رویه‌ای دارد، باید بتوان توضیح داد این اجزا با چه شبکه‌ها یا نواحی مغزی ارتباط دارند.

این نگاشت نباید ساده‌انگارانه باشد. مغز شبکه‌ای و تعاملی عمل می‌کند و هیچ مرز کاملاً تیزی میان همهٔ کارکردها وجود ندارد. اما نگاشت تقریبی برای اعتبار شناختی مدل مهم است. برای نمونه:
\begin{longtable}{>{\raggedleft\arraybackslash}p{0.34\textwidth}p{0.50\textwidth}}
\toprule
\textbf{جزء معماری} & \textbf{نگاشت تقریبی مغزی} \\
\midrule
ماژول بینایی & نواحی بینایی، به‌ویژه لوب پس‌سری \\
ماژول حرکتی & نواحی حرکتی در لوب پیشانی و مدارهای مرتبط \\
تصمیم‌گیری و کنترل شناختی & نواحی پیش‌پیشانی \\
حافظهٔ اعلانی & هیپوکامپ و ساختارهای لوب گیجگاهی میانی \\
حافظهٔ رویه‌ای & عقده‌های قاعده‌ای و مدارهای قشری-زیرقشری مرتبط \\
\bottomrule
\end{longtable}

این جدول یک نقشهٔ آموزشی است، نه ادعای جداسازی کامل. هرکدام از این کارکردها با شبکه‌های گسترده‌تر و تعامل میان چند ناحیه انجام می‌شوند.

\subsection{\eng{ACT-R}}
\term{ACT-R}{Adaptive Control of Thought-Rational} یکی از معماری‌های شناختی اثرگذار است که عمدتاً با کارهای \eng{John R. Anderson} و همکارانش شناخته می‌شود. ریشه‌های تاریخی آن به نظریه‌های \eng{ACT} در دههٔ ۱۹۷۰ برمی‌گردد، اما خود \eng{ACT-R} در نسخه‌های بعدی توسعه یافته و همچنان در مدل‌سازی شناختی استفاده می‌شود.

ایدهٔ کلی \eng{ACT-R} این است که شناخت انسان را می‌توان با چند جزء اصلی مدل کرد: ماژول‌ها، بافرها، حافظهٔ اعلانی، حافظهٔ رویه‌ای و سازوکار تطبیق الگو. در این معماری، دانش اعلانی معمولاً به شکل \eng{chunk}ها و دانش رویه‌ای به شکل \eng{production rule}ها نمایش داده می‌شود. وضعیت لحظه‌ای سیستم در محتوای بافرها دیده می‌شود و \eng{pattern matcher} قاعده‌ای را پیدا می‌کند که با وضعیت فعلی سازگار است.

در مثال جارو رباتیک، مسیر ساده‌شده چنین است:
\[
\begin{aligned}
\text{\lr{Environment}}
&\rightarrow
\text{\lr{Visual Module}}
\rightarrow
\text{\lr{Buffer}}\\
&\rightarrow
\text{\lr{Memory Retrieval}}
\rightarrow
\text{\lr{Production}}
\rightarrow
\text{\lr{Motor Module}}.
\end{aligned}
\]

\coursefigure{0.52}{act-r-architecture.png}{نمای شماتیک معماری \eng{ACT-R}: محیط از مسیر ماژول‌های بینایی و حرکتی با بافرها، حافظهٔ اعلانی، حافظهٔ رویه‌ای، تطبیق الگو و اجرای \eng{production} در ارتباط است.}

\subsubsection{حافظهٔ اعلانی}
حافظهٔ اعلانی در \eng{ACT-R} محل دانش صریح و قابل بیان است. در مثال ربات، وقتی بافر شامل «فنجان» و «آب» است، حافظهٔ اعلانی می‌تواند دانشی مانند این را بازیابی کند:
\begin{itemize}
  \item فنجان می‌تواند شکستنی باشد.
  \item آب یک مایع است.
  \item برخورد با فنجان می‌تواند باعث ریختن آب شود.
  \item ریختن آب زمین را خیس می‌کند.
\end{itemize}

این نوع دانش برای توضیح‌پذیری نیز مهم است. اگر سیستم مسیرش را عوض کند، می‌توان گفت: «فنجان پر از آب در مسیر بود و برخورد با آن می‌توانست باعث ریختن آب شود.» این بخش به بحث \eng{interpretability} در هوش مصنوعی نزدیک می‌شود، چون تصمیم مدل به دانش قابل بیان متصل می‌شود.

\subsubsection{حافظهٔ رویه‌ای}
حافظهٔ رویه‌ای در \eng{ACT-R} شامل \eng{production}ها یا قواعدی از جنس اگر-آنگاه است. این قواعد شبیه مهارت‌هایی‌اند که می‌گویند در یک وضعیت خاص چه عملی مناسب است:
\[
\text{اگر فنجان پر از آب در مسیر بود، مسیر را تغییر بده.}
\]
یا:
\[
\text{اگر مانع شکستنی در مسیر بود، سرعت را کم کن و آرام جابه‌جایش کن.}
\]

در انسان، بسیاری از مهارت‌ها پس از تمرین زیاد سریع و کم‌هزینه اجرا می‌شوند. حافظهٔ رویه‌ای نیز دقیقاً به همین جنبهٔ «چگونه انجام دادن» مربوط است، نه فقط «چه چیزی دانستن».

\subsubsection{\eng{Pattern Matching}}
\eng{Pattern Matching} مرحله‌ای است که در آن وضعیت فعلی بافرها و اطلاعات بازیابی‌شده با قواعد حافظهٔ رویه‌ای مقایسه می‌شود. اگر وضعیت فعلی با شرط یک \eng{production} سازگار باشد، آن قاعده نامزد اجرا می‌شود.

برای مثال:
\begin{itemize}
  \item بافر: فنجان پر از آب در مسیر است.
  \item حافظهٔ اعلانی: فنجان شکستنی است و آب ممکن است بریزد.
  \item حافظهٔ رویه‌ای: اگر مانع پرخطر در مسیر بود، با آن مثل مانع معمولی برخورد نکن.
  \item نتیجه: قاعدهٔ تغییر مسیر یا جابه‌جایی آرام فعال می‌شود.
\end{itemize}

در \eng{ACT-R}، شناخت به‌صورت توالی‌ای از فعال‌شدن و اجرای \eng{production}ها پیش می‌رود. بنابراین فرایند فقط یک خروجی نهایی نیست؛ زنجیره‌ای از وضعیت‌ها و تغییر وضعیت‌هاست.

\subsubsection{\eng{Production Execution}}
پس از انتخاب قاعده، \eng{production} اجرا می‌شود. خروجی این مرحله باید بتواند وضعیت بافرها را تغییر دهد و در نهایت به دستور حرکتی تبدیل شود. برای ربات، «مسیر را تغییر بده» باید به دستورهایی مانند کاهش سرعت، چرخش چرخ‌ها یا فعال‌کردن بازو تبدیل شود.

نکتهٔ مهم این است که خروجی اجرای \eng{production} باید به حافظهٔ کوتاه‌مدت یا بافر برگردد. سیستم باید بداند چه تصمیمی گرفته و چه عملی را شروع کرده است. این بازگشت باعث می‌شود توالی عمل‌ها حفظ شود، عمل بعدی با عمل قبلی سازگار باشد و رفتار سیستم از هم گسسته نشود.

\subsection{اعتبارسنجی معماری شناختی}
اعتبارسنجی معماری شناختی فقط با مقایسهٔ خروجی انجام نمی‌شود. ممکن است انسان و مدل هر دو به پاسخ درست برسند، اما مسیر رسیدن آن‌ها کاملاً متفاوت باشد. در علوم شناختی، خود مسیر اهمیت دارد.

برای اعتبارسنجی، می‌توان یک \eng{task} مشخص طراحی کرد و آن را هم به انسان و هم به مدل داد. در انسان، ابزارهایی مانند \eng{fMRI}، \eng{EEG}، زمان واکنش یا \eng{eye-tracking} می‌توانند اطلاعاتی دربارهٔ توالی و زمان‌بندی فعالیت‌ها بدهند. در مدل نیز می‌توان دید کدام ماژول در هر مرحله فعال شده است.

اگر مدل ابتدا ماژول بینایی را فعال کند، سپس وضعیت را در بافر قرار دهد، بعد حافظهٔ اعلانی و رویه‌ای را درگیر کند و در نهایت عمل حرکتی تولید کند، باید بررسی شود آیا دادهٔ انسانی نیز جریان مشابهی را نشان می‌دهد یا نه. هرچه توالی فعال‌شدن اجزا، زمان‌بندی و الگوی خطا با انسان سازگارتر باشد، معماری از نظر شناختی معتبرتر است.

\section{حافظه در انسان و مدل}

\subsection{حافظهٔ حسی، کوتاه‌مدت و کاری}
حافظه در معماری شناختی فقط یک محل ذخیره‌سازی ساده نیست. اگر قرار است مدل مهارت شناختی داشته باشد، باید روشن شود چه نوع اطلاعاتی را موقت نگه می‌دارد، چه دانشی را در بلندمدت ذخیره می‌کند و چگونه میان این دو سطح رفت‌وبرگشت انجام می‌دهد.

حافظه را می‌توان ابتدا به‌صورت ساده به حافظهٔ کوتاه‌مدت و بلندمدت تقسیم کرد. حافظهٔ کوتاه‌مدت شبیه یک بافر یا \eng{cache} محدود است. اطلاعات تازه ابتدا وارد این فضا می‌شوند و ممکن است خیلی زود از بین بروند.

مثلاً اگر کسی آدرس اتاقی را برای لحظه‌ای بشنود، ممکن است همان روز از آن استفاده کند، اما اگر ماه بعد دوباره به آن نیاز داشته باشد، احتمالاً باید دوباره بپرسد. این مثال نشان می‌دهد حضور موقت یک اطلاعات در ذهن به معنی ورود پایدار آن به حافظهٔ بلندمدت نیست.

در ادبیات جدید، باید میان \term{حافظهٔ کوتاه‌مدت}{Short-Term Memory} و \term{حافظهٔ کاری}{Working Memory} تمایز گذاشت. حافظهٔ کوتاه‌مدت بیشتر بر نگهداری موقت تأکید دارد، در حالی که حافظهٔ کاری علاوه بر نگهداری، پردازش فعال اطلاعات فعلی را نیز شامل می‌شود. این دو مفهوم گاهی نزدیک به هم استفاده می‌شوند، اما از نظر نظری یکسان نیستند.

\subsection{حافظهٔ بلندمدت}
اطلاعاتی که با تمرین، تکرار، استفاده یا تجربهٔ مهم تثبیت می‌شوند، می‌توانند وارد حافظهٔ بلندمدت شوند. حافظهٔ بلندمدت یک مخزن یکنواخت نیست و انواع مختلفی دارد. دو دستهٔ اصلی در اینجا حافظهٔ اعلانی و حافظهٔ رویه‌ای هستند.

\subsubsection{حافظهٔ اعلانی}
\term{حافظهٔ اعلانی}{Declarative Memory} مربوط به دانشی است که می‌توان آن را به‌صورت صریح بیان کرد: واقعیت‌ها، مفاهیم، رویدادها و اطلاعاتی که می‌توان درباره‌شان گزارش داد. مثال‌ها:
\begin{itemize}
  \item دانستن اینکه آب زمین را خیس می‌کند.
  \item دانستن اینکه فنجان ممکن است شکستنی باشد.
  \item دانستن آدرس یک مکان.
  \item دانستن یک تعریف یا مفهوم علمی.
\end{itemize}

از نظر مغزی، حافظهٔ اعلانی به‌ویژه با هیپوکامپ و ساختارهای لوب گیجگاهی میانی ارتباط دارد، هرچند ذخیره و بازیابی دانش در شبکه‌های گسترده‌تری انجام می‌شود. بنابراین بهتر است نگوییم همهٔ حافظهٔ اعلانی «در هیپوکامپ ذخیره می‌شود»؛ دقیق‌تر این است که هیپوکامپ و نواحی پیرامونی برای شکل‌گیری و بازیابی بسیاری از خاطرات اعلانی نقش کلیدی دارند.

در مقایسهٔ آموزشی با مدل‌های زبانی، می‌توان گفت مدل‌هایی مانند \eng{ChatGPT} در بازیابی و ترکیب دانش زبانی و واقعیت‌های آموخته‌شده قوی به نظر می‌رسند؛ بنابراین از این جهت به حافظهٔ اعلانی تشبیه می‌شوند. این تشبیه کامل نیست، چون مدل زبانی حافظهٔ انسانی ندارد، اما برای فهم تفاوت «دانستن» و «انجام دادن» مفید است.

\coursefigure{0.50}{hippocampus-and-declarative-memory.png}{هیپوکمپ و ساختارهای پیرامونی آن در تثبیت و بازیابی حافظهٔ اعلانی نقش کلیدی دارند؛ این تصویر برای پیوند بحث حافظهٔ صریح با نگاشت مغزی آمده است.}

\subsubsection{حافظهٔ رویه‌ای}
\term{حافظهٔ رویه‌ای}{Procedural Memory} مربوط به مهارت‌ها، عادت‌ها و روش‌های انجام کار است. مثال‌ها:
\begin{itemize}
  \item راه رفتن
  \item رانندگی
  \item آشپزی
  \item ترمز گرفتن
  \item حرکت دادن یک ابزار یا بازوی ربات
\end{itemize}

حافظهٔ رویه‌ای بیشتر با «چگونه انجام دادن» سروکار دارد. ممکن است فردی نتواند به یاد بیاورد دیشب چه غذایی خورده، اما همان لحظه بتواند آشپزی کند. این جدایی نسبی نشان می‌دهد دانستن یک واقعیت و اجرای یک مهارت یکسان نیستند.

از نظر مغزی، حافظهٔ رویه‌ای با عقده‌های قاعده‌ای، به‌ویژه مدارهای مرتبط با \eng{striatum}، و نیز مدارهای حرکتی و یادگیری عادت ارتباط دارد. در مثال هوش مصنوعی، خودروی خودران می‌تواند تا حدی به حافظهٔ رویه‌ای تشبیه شود، چون بخشی از مهارت رانندگی را اجرا می‌کند. البته خودروی خودران هم دانش اعلانی، ادراک و تصمیم‌گیری دارد؛ پس این فقط یک تشبیه آموزشی است.

\coursefigure{0.50}{basal-ganglia-and-procedural-memory.png}{عقده‌های قاعده‌ای در یادگیری مهارت، عادت و حافظهٔ رویه‌ای نقش دارند؛ این نگاشت کمک می‌کند تفاوت «دانستن» و «انجام دادن» از نظر مغزی روشن‌تر شود.}

\begin{correctionbox}{دقت دربارهٔ \eng{Chain of Thought}}
\eng{Chain of Thought} را نباید مستقیماً معادل حافظهٔ رویه‌ای دانست. این روش بیشتر به آشکارسازی یا هدایت گام‌های استدلال کمک می‌کند و به بازیابی دانش و استدلال صریح نزدیک‌تر است. ممکن است در اجرای یک روش مرحله‌به‌مرحله اثر بگذارد، اما خودش مدل کامل حافظهٔ رویه‌ای انسان نیست.
\end{correctionbox}

\subsection{انتقال اطلاعات از حافظهٔ کوتاه‌مدت به بلندمدت}
اطلاعات معمولاً با تمرین، تکرار و استفاده وارد حافظهٔ بلندمدت می‌شوند. اما برخی شرایط می‌توانند این انتقال را بسیار سریع‌تر یا قوی‌تر کنند.

\subsubsection{احساس شدید}
اگر فرد هنگام دریافت یک اطلاعات یا تجربه، احساس شدیدی داشته باشد، احتمال ماندگاری آن بیشتر می‌شود. تصادف، تولد فرزند، یک خبر بسیار خوشحال‌کننده یا یک تجربهٔ ترسناک می‌تواند همراه با جزئیات زیادی در حافظه بماند: زمان، مکان، فصل، وضعیت هوا یا حتی صحنه‌های فرعی.

\subsubsection{تعجب و غیرمنتظره بودن}
چیزی که فرد را غافلگیر کند، راحت‌تر در حافظه می‌ماند. به همین دلیل در آموزش، تجربهٔ غیرمنتظره، آزمایش عملی یا مثال متفاوت می‌تواند یادگیری را تقویت کند. البته غافلگیری باید به فهم موضوع وصل شود؛ صرفاً عجیب بودن کافی نیست.

\subsubsection{تمایز و برجستگی}
اطلاعاتی که از زمینهٔ معمول جدا می‌شوند، شانس بیشتری برای ماندن دارند. اگر فردی ظاهر یا رفتاری بسیار متمایز داشته باشد، احتمال دارد نام یا تصویر او بهتر در ذهن بماند. در آموزش نیز برجسته‌سازی نکتهٔ مهم می‌تواند به ماندگاری کمک کند.

\subsubsection{علامت‌گذاری و ارتباط‌سازی}
روش‌هایی مانند علامت‌گذاری، ساخت مثال، تصویرسازی ذهنی، تجسم‌سازی و پیدا کردن ارتباط میان مفاهیم کمک می‌کنند اطلاعات جدید با ساختارهای قبلی ذهن پیوند بخورد. این پیوندها مسیر بازیابی را تقویت می‌کنند و احتمال انتقال به حافظهٔ بلندمدت را بالا می‌برند.

از منظر آموزشی، این پرسش خودش یک خط پژوهشی مهم است: چگونه می‌توان یادگیری را طوری طراحی کرد که اطلاعات مهم سریع‌تر، عمیق‌تر و پایدارتر وارد حافظهٔ بلندمدت شوند؟

\subsection{اختلالات و شواهد شناختی}
تقسیم‌بندی حافظه فقط زمانی ارزش نظری دارد که بتواند داده‌های انسانی، از جمله اختلالات، را توضیح دهد. اگر می‌گوییم حافظهٔ کوتاه‌مدت، بلندمدت، اعلانی و رویه‌ای از هم متمایزند، باید بتوانیم نشان دهیم در برخی شرایط یکی از این بخش‌ها بیشتر آسیب می‌بیند و دیگری بهتر باقی می‌ماند.

\subsubsection{استرس و بازیابی حافظه}
استرس می‌تواند بازیابی اطلاعات اعلانی را مختل کند. فرد ممکن است چیزی را بداند، اما در لحظهٔ فشار نتواند آن را توضیح دهد یا به یاد آورد. این حالت را می‌توان به اختلال در بازیابی \eng{fact}ها یا توضیح‌های صریح ربط داد.

اگر استرس روی خود عمل صحبت‌کردن، لکنت، کنترل تنفس یا اجرای مهارت اثر بگذارد، بخشی از مسئله می‌تواند با جنبه‌های رویه‌ای نیز مرتبط شود. بنابراین بهتر است استرس را فقط به یک نوع حافظه محدود نکنیم؛ اثر آن به نوع وظیفه و سطح درگیری سیستم بستگی دارد.

\subsubsection{پارکینسون و حافظهٔ رویه‌ای}
در بیماری پارکینسون، به‌ویژه در مراحل اولیه، اختلالات حرکتی و مهارت‌های رویه‌ای برجسته‌اند. لرزش، کندی حرکت، دشواری در شروع حرکت و تغییر در راه رفتن می‌تواند باعث شود فرد ابتدا تصور کند مشکل فقط از پا، کمر یا عضله است؛ در حالی که بخش مهمی از مسئله به مدارهای عصبی کنترل حرکت و یادگیری رویه‌ای مربوط است.

این مثال با نقش عقده‌های قاعده‌ای در حرکت، عادت و یادگیری رویه‌ای سازگار است. البته پارکینسون فقط یک اختلال سادهٔ حافظهٔ رویه‌ای نیست و در طول زمان می‌تواند جنبه‌های شناختی دیگری را نیز درگیر کند. بنابراین باید آن را به‌عنوان شاهدی برای تعامل سیستم‌های حرکتی، رویه‌ای و شناختی فهمید، نه به‌عنوان جداسازی کامل یک ماژول از بقیه.

\subsection{حافظه در مدل‌های زبانی}
در این چارچوب، مفاهیم حافظهٔ شناختی می‌توانند برای فکر کردن دربارهٔ مدل‌های زبانی مفید باشند. این نگاشت استعاری است و نباید آن را یکی‌گرفتن مستقیم مدل زبانی با مغز انسان دانست.

\subsubsection{پارامترها به‌عنوان حافظهٔ بلندمدت}
وزن‌ها و پارامترهای آموزش‌دیدهٔ مدل را می‌توان به حافظهٔ بلندمدت تشبیه کرد. دانشی که مدل در طول آموزش از داده‌ها جذب کرده، در این پارامترها توزیع شده است. وقتی مدل به پرسشی پاسخ می‌دهد، بخشی از این دانش آموخته‌شده را فعال و ترکیب می‌کند.

\subsubsection{\eng{Prompt} و \eng{Context Window} به‌عنوان حافظهٔ کاری}
\eng{Prompt} و \eng{context window} را می‌توان به حافظهٔ کاری یا بافر فعلی مدل تشبیه کرد. اطلاعاتی که در متن ورودی قرار دارد، در همان تعامل جاری در دسترس مدل است. اگر متن طولانی، پرنویز یا متناقض شود، مدل ممکن است در نگهداری و استفادهٔ درست از اطلاعات مهم دچار مشکل شود؛ مسئله‌ای که در فصل‌های بعد با عنوان بار شناختی و محدودیت حافظهٔ کاری بررسی می‌شود.

\subsubsection{\eng{RAG} به‌عنوان حافظهٔ بیرونی}
\eng{Retrieval-Augmented Generation} یا \eng{RAG} را می‌توان شبیه حافظهٔ بیرونی دانست. مدل به‌جای تکیهٔ کامل بر پارامترهای خودش، اطلاعات مرتبط را از منبعی بیرونی بازیابی می‌کند و سپس بر اساس آن پاسخ می‌دهد. این سازوکار به‌ویژه وقتی مهم است که دانش لازم بسیار خاص، تازه، حجیم یا نیازمند ارجاع دقیق باشد.

\section{نظریهٔ بار شناختی و حافظهٔ کاری}
تا اینجا حافظه به‌عنوان یکی از ماژول‌های اصلی معماری شناختی معرفی شد. در این فصل، تمرکز از «انواع حافظه» به محدودیت‌های عملی حافظهٔ کاری منتقل می‌شود. پرسش اصلی این است: وقتی اطلاعات زیاد، پراکنده یا مزاحم باشند، انسان و مدل چگونه افت می‌کنند؟

\subsection{\eng{Cognitive Load Theory}}
\term{نظریهٔ بار شناختی}{Cognitive Load Theory} بر این ایده تکیه دارد که حافظهٔ کاری انسان محدود است. یادگیری، حل مسئله و استدلال وقتی دشوار می‌شوند که مقدار اطلاعاتی که باید هم‌زمان نگه داشته و پردازش شود از ظرفیت مؤثر حافظهٔ کاری فراتر برود.

این نظریه در آموزش اهمیت ویژه‌ای دارد، چون طراحی بدِ محتوا می‌تواند بار اضافه ایجاد کند. اگر اسلاید، تمرین، منو، رابط کاربری یا توضیح آموزشی هم‌زمان اطلاعات زیادی را بدون ساختار مناسب ارائه کند، بخشی از توان شناختی خواننده صرف مدیریت آشفتگی می‌شود، نه فهم مفهوم اصلی.

در ادبیات کلاسیک آموزش، معمولاً میان چند نوع بار شناختی تمایز گذاشته می‌شود: دشواری ذاتی خود مطلب، بار اضافی ناشی از شیوهٔ ارائه، و باری که صرف ساختن طرحوارهٔ مفهومی مفید می‌شود. در این منبع، تمرکز اصلی روی ایدهٔ عمومی‌تر است: حافظهٔ کاری محدود است و طراحی محاسباتی یا آموزشی باید این محدودیت را جدی بگیرد.

\subsection{ظرفیت محدود حافظهٔ کاری}
حافظهٔ کاری فقط محل نگهداری موقت نیست؛ محل پردازش فعال اطلاعات فعلی نیز هست. وقتی شماره تلفنی را برای چند ثانیه در ذهن نگه می‌داریم تا آن را وارد گوشی کنیم، از حافظهٔ کاری استفاده کرده‌ایم. اگر هم‌زمان کسی با ما حرف بزند، عددها طولانی باشند یا الگوی آشنایی نداشته باشند، خطا بیشتر می‌شود.

برای بررسی ظرفیت حافظهٔ کاری می‌توان آزمایش‌های گوناگون طراحی کرد: نمایش سریع چند تصویر و پرسیدن تعداد یا نوع آیتم‌ها، پخش چند صدا و سنجش تشخیص آن‌ها، ارائهٔ رشته‌های عددی با طول‌های مختلف، یا دادن چند مزه و درخواست دسته‌بندی. وجه مشترک این آزمایش‌ها این است که با افزایش تعداد یا پیچیدگی آیتم‌ها، عملکرد افت می‌کند.

همین ایده برای مدل‌های زبانی نیز قابل ترجمه است. به جای نمایش تصویر یا رشتهٔ عدد به انسان، می‌توان \eng{prompt}هایی طراحی کرد که اطلاعات طولانی، پراکنده، چندمرحله‌ای یا همراه با نویز دارند و سپس دید آیا مدل همچنان می‌تواند اطلاعات مرتبط را نگه دارد و از آن‌ها استفاده کند یا نه.

\subsection{\texorpdfstring{عدد \eng{7 \(\pm\) 2}}{عدد 7 ± 2}}
مقالهٔ مشهور جورج میلر با عنوان \eng{The Magical Number Seven, Plus or Minus Two} در سال ۱۹۵۶ یکی از نقاط کلاسیک بحث دربارهٔ ظرفیت حافظهٔ کوتاه‌مدت است. برداشت رایج از این مقاله این است که انسان می‌تواند حدود \(7 \pm 2\) واحد را در حافظهٔ کوتاه‌مدت نگه دارد.

اما این عدد نباید به‌عنوان قانون ثابت همهٔ انسان‌ها و همهٔ موقعیت‌ها فهمیده شود. پژوهش‌های بعدی نشان داده‌اند ظرفیت مؤثر به نوع ماده، طول کلمات، آشنایی فرد با محتوا، روش اندازه‌گیری، سن، تمرین و راهبردهای سازمان‌دهی بستگی دارد. پیام اصلی برای این بحث، خود عدد هفت نیست؛ پیام اصلی این است که حافظهٔ کاری محدود است و واحد مهم در بسیاری از موقعیت‌ها \eng{chunk} است.

شکل زیر صفحهٔ آغازین مقالهٔ میلر را نشان می‌دهد تا روشن باشد این بحث از یک مشاهدهٔ تجربی دربارهٔ ظرفیت پردازش اطلاعات شروع شده و در طول زمان به نظریه‌های دقیق‌تر حافظهٔ کاری و بار شناختی رسیده است.

\coursefigure{0.74}{miller-1956-cognitive-capacity-paper.png}{صفحهٔ آغازین مقالهٔ کلاسیک جورج میلر دربارهٔ محدودیت ظرفیت پردازش اطلاعات؛ این مقاله نقطهٔ تاریخی مهمی برای بحث حافظهٔ کوتاه‌مدت و مفهوم چانک است.}

\begin{correctionbox}{دقت دربارهٔ عدد میلر}
عدد \(7 \pm 2\) یک بیان کلاسیک و آموزشی مفید است، اما ظرفیت حافظهٔ کاری را نمی‌توان همیشه با یک عدد ثابت توضیح داد. در بسیاری از بحث‌های جدید، ظرفیت مؤثر کمتر و وابسته‌تر به نوع اطلاعات و راهبردهای فرد در نظر گرفته می‌شود.
\end{correctionbox}

\subsection{\eng{Chunking}}
\eng{Chunking} راهی برای کاهش بار حافظهٔ کاری است. انسان به جای نگه داشتن تعداد زیادی جزء جداگانه، آن‌ها را به واحدهای معنادار بزرگ‌تر تبدیل می‌کند. به این ترتیب، حافظهٔ کاری به جای چندین جزء پراکنده، چند واحد سازمان‌یافته را نگه می‌دارد.

\subsubsection{تعریف \eng{Chunk}}
\term{چانک}{Chunk} یعنی گروهی از اطلاعات کوچک‌تر که برای فرد به‌صورت یک واحد معنادار در می‌آیند. برای مثال:
\[
\text{\lr{1 3 6 9}}
\]
اگر این چهار رقم جداگانه دیده شوند، چهار آیتم‌اند. اما اگر فرد آن‌ها را سال \eng{1369} یا بخشی از یک الگوی آشنا بداند، کل رشته می‌تواند یک \eng{chunk} شود.

بنابراین اندازهٔ واقعی حافظهٔ کاری فقط به تعداد خام نشانه‌ها بستگی ندارد؛ به این بستگی دارد که ذهن چگونه آن نشانه‌ها را سازمان‌دهی می‌کند.

\subsubsection{نقش دانش قبلی در \eng{Chunking}}
\eng{Chunking} به حافظهٔ بلندمدت وابسته است. ما برای فشرده‌سازی اطلاعات جدید از دانش قبلی، الگوهای آشنا، معنا، تصویر ذهنی و تداعی‌های قبلی استفاده می‌کنیم. کسی که یک رشتهٔ عددی را به تاریخ تولد، شمارهٔ ماشین، وزن ورزشی یا الگوی ریاضی ربط می‌دهد، آن را بهتر از کسی نگه می‌دارد که فقط رقم‌ها را جداگانه حفظ می‌کند.

به همین دلیل تخصص می‌تواند ظرفیت مؤثر حافظهٔ کاری را تغییر دهد. متخصص شطرنج ممکن است چیدمان صفحه را به چند الگوی معنادار تبدیل کند، در حالی که فرد مبتدی همان صفحه را مجموعه‌ای از مهره‌های پراکنده می‌بیند.

\subsubsection{تفاوت‌های فردی}
تفاوت افراد در حافظهٔ کاری فقط به ظرفیت خام محدود نمی‌شود. افراد در نوع و کیفیت \eng{chunking} نیز تفاوت دارند. یک نفر چانک تصویری می‌سازد، دیگری چانک زبانی، دیگری از داستان یا تداعی استفاده می‌کند و فردی دیگر الگوی عددی یا فضایی پیدا می‌کند.

بنابراین نظریهٔ خوب باید بتواند توضیح دهد چرا دو نفر با دادهٔ ظاهراً یکسان عملکرد متفاوت دارند. تفاوت می‌تواند از اندازهٔ چانک‌ها، نوع چانک‌ها، توانایی اتصال اطلاعات جدید به حافظهٔ بلندمدت، یا راهبردهای سازمان‌دهی ذهنی ناشی شود.

\subsection{کاربردهای \eng{Chunking}}
\eng{Chunking} فقط یک مفهوم آزمایشگاهی نیست. در طراحی آموزشی، طراحی اسلاید، طراحی رابط کاربری، طراحی منوها، بازی‌ها، شماره تلفن‌ها و سازمان‌دهی اطلاعات پیچیده از همین اصل استفاده می‌شود. اگر کاربر یا خواننده با تعداد زیادی آیتم بی‌ساختار روبه‌رو شود، حافظهٔ کاری او سریع‌تر اشباع می‌شود. اما اگر همان اطلاعات در گروه‌های معنادار، مرحله‌های کوچک‌تر یا ساختارهای قابل پیش‌بینی ارائه شود، پردازش ساده‌تر می‌شود.

\subsection{\eng{Multi-Hop Reasoning} به‌عنوان آزمون حافظهٔ کاری}
\term{استدلال چندمرحله‌ای}{Multi-Hop Reasoning} وظیفه‌ای است که پاسخ نهایی در آن از یک گام ساده به دست نمی‌آید. مدل یا انسان باید چند قطعه اطلاعات را نگه دارد، رابطهٔ میان آن‌ها را بفهمد، نتیجهٔ میانی بسازد و از نتیجهٔ میانی برای گام بعدی استفاده کند.

در یک مثال ساده:
\begin{itemize}
  \item انسان‌ها فانی‌اند.
  \item سقراط انسان است.
  \item پس سقراط فانی است.
  \item اگر فانی‌ها معمولاً تا سن خاصی عمر کنند، نتیجهٔ بعدی نیز به این زنجیره وابسته می‌شود.
\end{itemize}

یا در پرسش خانوادگی «پدرِ همسرِ برادرِ مادر من چه کسی است؟» باید چند رابطهٔ پشت سر هم پردازش شود. هر گام باید تا گام بعدی در حافظهٔ کاری فعال بماند. همین ویژگی، \eng{multi-hop reasoning} را برای سنجش بار شناختی در انسان و مدل مناسب می‌کند.

\subsection{محدودیت حافظهٔ کاری در مدل‌های زبانی}
در مدل‌های زبانی، معادل آزمایش حافظهٔ کاری را می‌توان با طراحی \eng{prompt}های دشوار انجام داد. هدف این است که نشان دهیم:
\begin{enumerate}
  \item مدل نیز در استفاده از زمینهٔ فعلی محدودیت دارد.
  \item با افزایش بار اطلاعاتی، عملکرد افت می‌کند.
  \item نویز، پراکندگی شواهد و تغییر موضوع می‌توانند بازیابی اطلاعات مرتبط را مختل کنند.
  \item سپس می‌توان با الهام از انسان، راهکارهایی مانند \eng{chunking}، خلاصه‌سازی، تقسیم مسئله یا مدیریت حافظه را بررسی کرد.
\end{enumerate}

در سال‌های اخیر، برخی کارها همین ایده را با عنوان بار شناختی محاسباتی در مدل‌های زبانی بررسی کرده‌اند. یکی از مسیرهای مطرح، سنجش مدل‌ها در \eng{multi-hop reasoning} همراه با دست‌کاری مقدار نویز یا ترتیب وظایف است.

\subsubsection{\eng{Context Saturation}}
\eng{Context Saturation} زمانی رخ می‌دهد که اطلاعات مهم در میان حجم زیادی از اطلاعات نامرتبط یا کم‌ارتباط پخش شود. در این حالت، مسئله فقط طول \eng{context window} نیست؛ مسئله این است که مدل باید از میان مقدار زیادی نویز، شواهد اصلی را پیدا کند.

ساختار آزمایش می‌تواند چنین باشد: ابتدا یک مسئلهٔ چندمرحله‌ای داده می‌شود، سپس مقدار زیادی متن نامرتبط یا اطلاعات مزاحم میان شواهد مرتبط اضافه می‌شود، و در نهایت عملکرد مدل اندازه‌گیری می‌شود. اگر با افزایش نویز، نرخ موفقیت کاهش یابد، می‌توان گفت زمینهٔ مدل از نظر شناختی اشباع شده است.

\subsubsection{\eng{Attentional Residue}}
\eng{Attentional Residue} به باقی‌ماندن اثر وظیفه یا موضوع قبلی روی وظیفهٔ فعلی اشاره دارد. در تجربهٔ عملی با مدل‌های زبانی، گاهی مدل سؤال جدید را با زمینهٔ قبلی قاطی می‌کند، یا هنوز در چارچوب مسئلهٔ قبلی پاسخ می‌دهد. در چنین حالتی، پاک‌کردن مکالمه، بازنویسی \eng{prompt} یا جداکردن وظایف می‌تواند کمک کند.

در طراحی آزمایش، می‌توان ابتدا چند بخش نامرتبط یا یک وظیفهٔ مزاحم به مدل داد، سپس مسئلهٔ اصلی را مطرح کرد و دید آیا مدل می‌تواند توجه خود را از اثر باقی‌ماندهٔ قبلی جدا کند یا نه.

\subsubsection{تفاوت دو نوع فشار}
هر دو پدیده با نویز و مزاحمت سروکار دارند، اما چیدمان آن‌ها متفاوت است:
\begin{itemize}
  \item در \eng{Context Saturation}، اطلاعات مرتبط در میان نویز پخش می‌شود.
  \item در \eng{Attentional Residue}، وظیفه یا نویز قبلی روی وظیفهٔ بعدی سایه می‌اندازد.
\end{itemize}
در هر دو حالت، افت عملکرد مدل نشانه‌ای از محدودیت در مدیریت زمینه، توجه و حافظهٔ کاری است.

\subsection{طراحی آزمایش برای سنجش بار شناختی مدل}
برای نشان دادن محدودیت حافظهٔ کاری در مدل زبانی، فقط طولانی کردن \eng{prompt} کافی نیست. باید طراحی آزمایش طوری باشد که عامل‌های مزاحم قابل کنترل و قابل تفسیر باشند. چند روش مطرح‌شده در این بحث:
\begin{itemize}
  \item افزایش طول \eng{prompt}
  \item اضافه کردن اطلاعات زیاد اما نامرتبط
  \item پخش کردن شواهد مرتبط در بخش‌های دور از هم
  \item اضافه کردن نویز میان اطلاعات مهم
  \item ارجاع دادن به بخش‌های قبلی مکالمه
  \item ایجاد تداخل میان موضوع قبلی و موضوع فعلی
  \item دادن یک کد یا دستور و سپس تغییر خواسته در ادامهٔ \eng{prompt}
\end{itemize}

در هر طراحی، باید روشن باشد چه چیزی اندازه‌گیری می‌شود: بازیابی اطلاعات مرتبط، حفظ نتیجهٔ میانی، مقاومت در برابر نویز، یا توانایی تغییر توجه از وظیفهٔ قبلی به وظیفهٔ جدید.

\subsection{تکامل آزمون‌های هوش مصنوعی}
این بحث در ادامهٔ تکامل ارزیابی هوش مصنوعی قرار می‌گیرد. در ابتدا آزمون‌هایی مانند \eng{Turing Test} برجسته بودند. بعدتر در یادگیری ماشین، معیارهایی مانند \eng{accuracy}، \eng{precision} و \eng{recall} روی \eng{test set} اهمیت پیدا کردند. با پیشرفت مدل‌ها، مسئلهٔ \term{تعمیم خارج از توزیع}{Out-of-Distribution Generalization} مطرح شد.

اکنون یکی از مسیرهای مهم، آزمون‌های شناختی است. به جای اینکه فقط بپرسیم مدل روی یک بنچمارک چند درصد امتیاز می‌گیرد، می‌پرسیم: آیا مدل مانند انسان محدودیت حافظهٔ کاری دارد؟ آیا خطاهای مدل با خطاهای انسان همبستگی دارد؟ آیا در مهارت‌هایی مانند استدلال، حافظه، برنامه‌ریزی و انتزاع ضعف نظام‌مند دارد؟ آیا می‌توان با الهام از نظریه‌های شناختی این ضعف را کاهش داد؟

\subsection{راهکارهای شناختی}
وقتی محدودیت حافظهٔ کاری شناسایی شد، مرحلهٔ بعد طراحی راهکار است. در انسان، راهکارهایی مانند \eng{chunking}، تقسیم مسئله، ساختن نشانه‌های بازیابی، خلاصه‌سازی و استفاده از دانش قبلی به کاهش بار شناختی کمک می‌کنند. در مدل‌های زبانی نیز می‌توان از ایده‌های مشابه استفاده کرد: تقسیم مسئله به زیرمسئله، خلاصه کردن شواهد، ساختن حافظهٔ بیرونی، یا طراحی \eng{prompt}هایی که اطلاعات را مرحله‌به‌مرحله سازمان‌دهی می‌کنند.

\subsubsection{\eng{Divide and Conquer}}
\eng{Divide and Conquer} یعنی مسئلهٔ بزرگ را به چند مسئلهٔ کوچک‌تر تقسیم کنیم. این ایده وقتی مفید است که نتیجهٔ هر بخش بتواند به‌صورت خلاصه و قابل استفاده برای گام بعدی نگه داشته شود. اگر مدل مجبور نباشد همهٔ جزئیات خام را هم‌زمان در \eng{context} نگه دارد، فشار حافظهٔ کاری کمتر می‌شود.

از نظر الگوریتمی، این ایده را می‌توان با رابطهٔ کلی زیر نمایش داد:
\[
T(n)=aT(n/b)+f(n)
\]
که در آن \(a\) تعداد زیرمسئله‌ها، \(n/b\) اندازهٔ هر زیرمسئله، و \(f(n)\) هزینهٔ شکستن مسئله و ترکیب پاسخ‌هاست. پیام آموزشی این فرمول این است که تقسیم درست مسئله می‌تواند هزینه را به‌شدت کم کند؛ برای مثال، در بعضی مسائل به جای هزینه‌های مرتبهٔ \(O(n^2)\) می‌توان به ساختارهایی نزدیک‌تر به \(O(n\log n)\) رسید. البته این مقایسه قانون عمومی همهٔ مسائل نیست و به ساختار دقیق مسئله و الگوریتم بستگی دارد.

یک مثال ساده، خانه‌تکانی است. اگر کل خانه یک‌باره موضوع کار باشد، همهٔ جزئیات هم‌زمان روی حافظهٔ کاری فشار می‌آورند. اما اگر خانه به اتاق‌ها، اتاق به کمد و کشو، و هر کشو به چند دستهٔ کوچک‌تر تقسیم شود، در هر لحظه فقط یک بخش فعال است. پس از تمام شدن هر بخش، لازم نیست همهٔ حرکت‌های کوچک در حافظه بماند؛ کافی است نتیجهٔ آن بخش به‌عنوان یک وضعیت خلاصه نگه داشته شود.

\subsubsection{\eng{Sub-goal} و مدیریت سلسله‌مراتبی حافظه}
در وظایف طولانی، هر \eng{sub-goal} می‌تواند نقش یک چانک سطح بالاتر را داشته باشد. مدل ابتدا یک زیرهدف را حل می‌کند، خروجی مهم آن را نگه می‌دارد، سپس به زیرهدف بعدی می‌رود. این ایده در \eng{agent}های بلندمدت مهم است، چون مدل باید بداند چه چیزی از مرحلهٔ قبلی را نگه دارد و چه جزئیاتی را می‌تواند کنار بگذارد.

در مقالهٔ \eng{HiAgent} این ایده به شکل صریح به مدیریت حافظهٔ کاری \eng{agent}های زبانی وصل می‌شود. در راهبرد ساده، همهٔ جفت‌های «عمل--مشاهده» در طول انجام کار وارد زمینهٔ مدل می‌شوند. این کار اطلاعات را کامل نگه می‌دارد، اما در وظایف بلندمدت زمینه را پر از جزئیات تکراری و کم‌اهمیت می‌کند. راه‌حل پیشنهادی این است که مدل ابتدا زیرهدف بسازد، عمل‌های مربوط به همان زیرهدف را دنبال کند، و وقتی زیرهدف کامل شد، جزئیات اجرایی آن را با یک مشاهدهٔ خلاصه جایگزین کند.

به این ترتیب، حافظهٔ کاری به‌جای یک فهرست بلند از همهٔ حرکت‌ها، چیزی شبیه این دارد:
\[
\left(g_1,s_1\right),\left(g_2,s_2\right),\ldots,\left(g_k,\text{جزئیات زیرهدف فعلی}\right)
\]
که در آن \(g_i\) زیرهدف و \(s_i\) خلاصهٔ وضعیت پس از انجام آن زیرهدف است. اگر بعداً جزئیات یک زیرهدف گذشته لازم شود، می‌توان آن بخش را دوباره بازیابی کرد؛ اما حالت پیش‌فرض این نیست که همهٔ جزئیات همیشه در حافظهٔ فعال بمانند.

\subsubsection{مثال \eng{Blocks World}}
در \eng{Blocks World} چند بلوک روی میز یا روی یکدیگر قرار دارند و هدف این است که به وضعیت نهایی مشخصی برسیم؛ مثلاً بلوک آبی روی بلوک نارنجی قرار گیرد. برای رسیدن به این هدف، مسئله را می‌توان به چند زیرهدف شکست:
\begin{enumerate}
  \item بلوک آبی \eng{clear} شود؛ یعنی چیزی روی آن نباشد.
  \item بلوک نارنجی \eng{clear} شود.
  \item بلوک آبی روی بلوک نارنجی گذاشته شود.
\end{enumerate}

نکتهٔ شناختی این نیست که زیرهدف در برنامه‌ریزی رباتیک چیز تازه‌ای است؛ زیرهدف‌ها از قبل در \eng{planning} وجود داشته‌اند. نکته این است که در اینجا زیرهدف به‌عنوان واحد حافظه نیز تفسیر می‌شود. پس از کامل شدن زیرهدف «بلوک نارنجی clear است»، لازم نیست همهٔ حرکت‌هایی که به این وضعیت منجر شده‌اند در حافظهٔ فعال بمانند. نتیجهٔ فشردهٔ زیرهدف برای ادامهٔ حل مسئله کافی است.

\subsubsection{مثال \eng{Gripper}}
در \eng{Gripper}، ربات باید اشیایی را بین اتاق‌ها جابه‌جا کند. در نتایج گزارش‌شدهٔ \eng{HiAgent}، چانک‌سازی سلسله‌مراتبی در بعضی وظایف مانند \eng{Blocksworld} سود آشکارتری داشت، اما در \eng{Gripper} بهبود موفقیت ایجاد نکرد و حتی شاخص پیشرفت کمی افت کرد، هرچند مصرف زمینه کمتر شد. تفسیر آموزشی این است که \eng{chunking} وقتی بیشترین فایده را دارد که مسئله واقعاً ساختار سلسله‌مراتبی و وابستگی میان زیرهدف‌ها داشته باشد. اگر کارها بسیار مستقل و خرد باشند، تقسیم اضافه ممکن است فقط هزینهٔ مدیریتی بسازد.

\subsubsection{محدودیت‌های \eng{Chunking}}
\eng{Chunking} همیشه راه‌حل کامل نیست. اگر مسئله ذاتاً ساختار تقسیم‌پذیر نداشته باشد، یا اگر تقسیم‌بندی غلط انجام شود، چانک‌ها می‌توانند اطلاعات مهم را پنهان کنند. همچنین در مدل‌های زبانی، ساختن چانک معمولاً به طراحی بیرونی، \eng{prompting} یا ماژول حافظه نیاز دارد؛ مدل همیشه خودش نمی‌داند چه چیزی را باید خلاصه کند و چه چیزی را باید نگه دارد.

تفاوت مهم انسان و مدل نیز همین‌جاست. انسان‌ها در بسیاری از موقعیت‌ها خودشان راهبرد می‌سازند: موضوع را دسته‌بندی می‌کنند، چیزهای کم‌اهمیت را کنار می‌گذارند، و تشخیص می‌دهند چه چیزی باید در حافظه بماند. در بسیاری از \eng{agent}های فعلی، ما هنوز باید صریحاً بگوییم زیرهدف چیست، چه چیزی چانک شود، چه چیزی فراموش شود و چه چیزی بازیابی شود. بنابراین راهبرد شناختی به مدل الهام می‌دهد، اما پیاده‌سازی آن هنوز به طراحی محاسباتی دقیق نیاز دارد.

\section{ترنسفورمر، توجه و حافظهٔ کاری}

\subsection{حافظه در \eng{Transformer}}
برای پیوند دادن بحث حافظهٔ کاری با معماری‌های امروزی، باید ببینیم در \eng{Transformer} چه چیزی را می‌توان به‌صورت تقریبی معادل حافظهٔ بلندمدت و کوتاه‌مدت دانست. این تشبیه کامل و زیستی نیست، اما برای تحلیل مدل مفید است.

حافظهٔ بلندمدت مدل در وزن‌ها و پارامترهای آن قرار دارد. دانشی که مدل در زمان آموزش از متن، گفتار، تصویر یا داده‌های دیگر آموخته، در همین پارامترها کدگذاری شده است. برای مثال، اینکه مدل واژهٔ \eng{data} را به چه حوزه‌هایی نزدیک می‌داند، یا چه رابطه‌هایی میان کلمات، ساختارهای نحوی و مفاهیم یاد گرفته است، محصول آموزش و پارامترهاست.

در مقابل، حافظهٔ کوتاه‌مدت یا حافظهٔ کاری به پردازش ورودی فعلی مربوط است: \eng{prompt}، جمله، متن، مکالمه یا دنباله‌ای که در زمان \eng{inference} پیش روی مدل قرار دارد. بخشی از معماری که بیشترین ارتباط را با این پردازش دارد، لایهٔ \eng{attention} است؛ زیرا مدل در آنجا تصمیم می‌گیرد هر توکن، برای ساختن بازنمایی فعلی خود، باید به کدام توکن‌های دیگر توجه کند.

\begin{correctionbox}{تشبیه، نه همسانی کامل}
وقتی می‌گوییم پارامترها شبیه حافظهٔ بلندمدت و \eng{attention} شبیه حافظهٔ کاری است، منظور یک نگاشت دقیق عصبی نیست. مدل زبانی مغز ندارد و حافظهٔ انسانی نیز صرفاً ماتریس وزن نیست. این تشبیه فقط کمک می‌کند نقش «دانش آموخته‌شده» و «پردازش زمینهٔ فعلی» را جدا کنیم.
\end{correctionbox}

\subsection{\eng{Contextual Embedding}}
مدل در ابتدا برای هر \eng{token} یک بازنمایی اولیه دارد. این بازنمایی از دانش آموخته‌شدهٔ مدل و جدول‌های \eng{embedding} یا لایه‌های ورودی می‌آید. اما هدف \eng{Transformer} این نیست که یک معنای ثابت برای هر کلمه نگه دارد؛ هدف این است که بازنمایی کلمه با توجه به زمینهٔ فعلی تغییر کند. به این بازنمایی وابسته به زمینه \term{بردار زمینه‌مند}{Contextual Embedding} می‌گوییم.

برای مثال، واژهٔ «شیر» در فارسی می‌تواند به حیوان، نوشیدنی یا وسیلهٔ باز و بسته کردن آب اشاره کند. اگر مدل فقط یک بردار ثابت برای «شیر» داشته باشد، نمی‌تواند این تفاوت‌ها را به‌خوبی نشان دهد. اما پس از لایه‌های \eng{self-attention}، بازنمایی «شیر» در «شیر در جنگل زندگی می‌کند» باید با بازنمایی آن در «شیر را در لیوان ریخت» متفاوت شود.

همین نکته برای واژهٔ \eng{data} نیز برقرار است. بسته به زمینه، \eng{data} ممکن است به پایگاه داده، یادگیری ماشین، مصورسازی، آمار، حریم خصوصی یا سیگنال آزمایشگاهی نزدیک شود. \eng{Contextual embedding} یعنی مدل معنای عملی توکن را در همان متن خاص بسازد، نه اینکه فقط به معنای میانگین و ثابت آن تکیه کند.

\subsubsection{حل ارجاع}
\term{حل ارجاع}{Coreference Resolution} یعنی مدل بفهمد چند عبارت مختلف به یک موجودیت واحد اشاره می‌کنند. در جملهٔ «علی رفت مدرسه. او کتابش را برداشت»، مدل باید تشخیص دهد «او» و «ـش» به علی برمی‌گردند. در سطح بازنمایی، انتظار داریم بردارهای مربوط به این ارجاع‌ها به اطلاعاتی دربارهٔ «علی» وصل شوند، نه اینکه ضمیرها جدا و بی‌ارتباط پردازش شوند.

این کار نمونهٔ روشنی از حافظهٔ کاری در پردازش زبان است. مدل باید موجودیتی را که در جملهٔ قبلی آمده نگه دارد، ضمیر را در جملهٔ بعدی ببیند، و میان آن‌ها پیوند بسازد.

\subsubsection{رفع ابهام معنای واژه}
\term{رفع ابهام معنای واژه}{Word Sense Disambiguation} یعنی تشخیص اینکه یک واژهٔ چندمعنا در متن فعلی کدام معنا را دارد. برای «شیر»، نشانه‌هایی مانند «جنگل»، «یال»، «لیوان»، «کلسیم»، «آب» یا «باز کردن» معنا را تعیین می‌کنند. مدل باید این نشانه‌ها را در زمینه پیدا کند و بازنمایی واژه را به معنای درست نزدیک کند.

پس \eng{contextual embedding} فقط یک اصطلاح فنی در \eng{NLP} نیست؛ همان نقطه‌ای است که مدل از حافظهٔ بلندمدت خود، یعنی دانش آماری و معنایی آموخته‌شده، برای پردازش وضعیت فعلی استفاده می‌کند.

\subsection{\eng{Query}، \eng{Key} و \eng{Value}}
در \eng{self-attention}، بازنمایی هر توکن با چند ماتریس وزنی به سه بردار تبدیل می‌شود:
\begin{itemize}
  \item \term{پرس‌وجو}{Query}: این توکن دنبال چه اطلاعاتی است؟
  \item \term{کلید}{Key}: این توکن چه نشانه‌ای برای یافتن خود ارائه می‌کند؟
  \item \term{مقدار}{Value}: اگر به این توکن توجه شود، چه اطلاعاتی منتقل می‌شود؟
\end{itemize}

اگر بردار اولیهٔ یک توکن مثلاً ابعاد \(1 \times 768\) داشته باشد، مدل با ضرب آن در ماتریس‌های آموخته‌شده، بردارهای \(Q\)، \(K\) و \(V\) می‌سازد. سپس \eng{query} یک توکن با \eng{key} توکن‌های دیگر مقایسه می‌شود. هرچه شباهت یا سازگاری آن‌ها بیشتر باشد، وزن توجه بیشتر می‌شود. در شکل ساده، وزن‌ها با ضرب داخلی و سپس \eng{softmax} ساخته می‌شوند:
\[
\operatorname{Attention}(Q,K,V)=\operatorname{softmax}\left(\frac{QK^{T}}{\sqrt{d_k}}\right)V
\]
که \(d_k\) بعد بردارهای کلید است. ضرب نهایی در \(V\) باعث می‌شود بازنمایی جدید هر توکن از ترکیب اطلاعات توکن‌هایی ساخته شود که برای آن مهم‌تر بوده‌اند.

در شکل زیر، همین تبدیل خطی به‌صورت تصویری دیده می‌شود: embeddingهای ورودی با وزن‌های آموخته‌شده ترکیب می‌شوند تا نمایش‌های \(Q\)، \(K\) و \(V\) ساخته شوند. این تصویر کمک می‌کند توجه را نه به‌عنوان یک مفهوم مبهم، بلکه به‌عنوان چند ضرب ماتریسی روی بازنمایی‌های فعلی ببینیم.

\coursefigure{0.88}{query-key-value-projections.png}{ساخت بردارهای \eng{Query}، \eng{Key} و \eng{Value} از embeddingهای ورودی با ضرب در وزن‌های آموخته‌شده و اضافه شدن بایاس.}

\subsection{\eng{Self-Attention} از دید حافظهٔ کاری}
\eng{Self-attention} را می‌توان به‌صورت آموزشی بخشی از سازوکار حافظهٔ کاری مدل دانست، چون روی ورودی فعلی کار می‌کند. پیش از \eng{attention}، مدل بازنمایی‌هایی دارد که عمدتاً از حافظهٔ بلندمدت و جایگاه توکن‌ها آمده‌اند. پس از \eng{attention}، هر توکن بازنمایی تازه‌ای دارد که به متن فعلی وابسته است.

برای مثال، در جمله‌ای که چند بار از «شیر» با معناهای متفاوت استفاده شده، \eng{attention} باید نشانه‌های اطراف هر نمونه را جمع کند تا بازنمایی مناسب بسازد. در مسئلهٔ چندمرحله‌ای، \eng{attention} باید میان قطعه‌های دور از هم پیوند برقرار کند. بنابراین اگر مدل در بازیابی شواهد مرتبط یا نادیده گرفتن اطلاعات نامرتبط ضعف داشته باشد، این ضعف را می‌توان تا حدی در نحوهٔ مدیریت توجه و زمینه جست‌وجو کرد.

البته در \eng{decoder-only transformer}های خودرگرسیو، هر توکن فقط به توکن‌های قبلی توجه می‌کند، چون تولید متن نباید به آینده دسترسی داشته باشد. در \eng{encoder-only}ها، مانند مدل‌های خانوادهٔ \eng{BERT}، توجه معمولاً دوسویه است و هر توکن می‌تواند از توکن‌های قبل و بعد استفاده کند. پس نوع ترنسفورمر بر شکل حافظهٔ کاری محاسباتی اثر می‌گذارد.

\coursefigure{0.76}{attention-weighted-values.png}{ضرب ماتریس توجه در بردارهای \eng{Value}: وزن‌های توجه مشخص می‌کنند هر توکن چه مقدار از اطلاعات توکن‌های دیگر را در بازنمایی خروجی خود دریافت کند.}

\subsection{آیا افزایش طول \eng{Context} کافی است؟}
یک پاسخ ساده به محدودیت حافظهٔ کاری مدل این است که پنجرهٔ زمینه را بزرگ‌تر کنیم. این کار در بسیاری از کاربردها مفید است، اما به‌تنهایی مسئله را حل نمی‌کند. انسان نیز اگر حافظهٔ خام بیشتری داشته باشد، بدون راهبرد سازمان‌دهی الزاماً بهتر فکر نمی‌کند. مسئلهٔ اصلی فقط «جا داشتن اطلاعات» نیست؛ مسئله این است که مدل بتواند اطلاعات مرتبط را بازیابی کند، اطلاعات نامرتبط را کنار بگذارد و مسیر استدلال را حفظ کند.

اگر \eng{prompt} طولانی باشد اما پر از نویز، تکرار یا موضوع‌های نامرتبط شود، افزایش طول زمینه حتی می‌تواند بار پردازشی بیشتری ایجاد کند. بنابراین هدف بهتر، افزایش بی‌قید حافظه نیست؛ هدف \term{مدیریت توجه و بازیابی}{Attention and Retrieval Management} است.

\subsection{کاهش یا حذف توجه به اطلاعات کم‌اهمیت}
یک ایدهٔ معماری این است که ارتباط‌هایی با \eng{attention score} پایین حذف یا صفر شوند تا مدل روی موارد مهم‌تر تمرکز کند. این ایده از نظر شهودی جذاب است، اما ساده نیست.

اول، وزن توجه نسبی است. ممکن است یک توکن نسبت به توکن دوم وزن کمی داشته باشد، اما نسبت به توکن سوم وزن بالایی داشته باشد. بنابراین «کم بودن» یک وزن در یک رابطه به‌تنهایی نشان نمی‌دهد آن توکن یا مسیر اطلاعاتی بی‌اهمیت است.

دوم، وابستگی‌ها می‌توانند زنجیره‌ای باشند. اگر \(A\) به \(B\) مربوط باشد و \(B\) به \(C\)، ممکن است \(A\) از راه \(B\) به \(C\) وابسته شود، حتی اگر توجه مستقیم \(A\) به \(C\) کوچک باشد. حذف خام وزن‌های کوچک می‌تواند این مسیرهای غیرمستقیم را خراب کند.

پس کاهش توجه باید با طراحی دقیق انجام شود: با پنجره‌های محلی، الگوهای پراکنده، بازیابی هدفمند، خلاصه‌سازی، یا لایه‌هایی که اجازه دهند اطلاعات مهم در سطح بالاتر دوباره ترکیب شود.

\subsection{\eng{Local Attention}}
\eng{Local Attention} یعنی هر توکن به جای مقایسه با همهٔ توکن‌های دنباله، فقط به یک پنجرهٔ محلی از توکن‌های نزدیک خود توجه کند. این کار از نظر محاسباتی مهم است، چون \eng{self-attention} کامل نسبت به طول دنباله هزینهٔ درجهٔ دوم دارد؛ اگر طول دنباله \(n\) باشد، مقایسهٔ همه با همه تقریباً از مرتبهٔ \(O(n^2)\) است.

با محلی کردن توجه، هزینه و مصرف حافظه کاهش می‌یابد و مدل مجبور نیست همهٔ ارتباط‌های ممکن را هم‌زمان بررسی کند. از نظر مفهومی، این ایده به \eng{divide and conquer} شباهت دارد: ورودی به بخش‌های کوچک‌تر تقسیم می‌شود، در هر بخش پردازش محلی انجام می‌گیرد، و سپس لایه‌های بعدی یا سازوکارهای مکمل باید اطلاعات بخش‌ها را با هم ترکیب کنند.

اما \eng{local attention} خطر مهمی دارد: وابستگی‌های دوربرد ممکن است از دست بروند. اگر پاسخ یک سؤال به جمله‌ای در ابتدای متن و جمله‌ای در انتهای متن وابسته باشد، پنجرهٔ محلی ساده ممکن است این دو را به هم وصل نکند. به همین دلیل، مدل‌های عملی معمولاً فقط به پنجرهٔ محلی خام تکیه نمی‌کنند؛ ممکن است از پنجره‌های لغزان، توکن‌های جهانی، الگوهای پراکنده یا لایه‌های سلسله‌مراتبی برای عبور اطلاعات بین بخش‌ها استفاده کنند.

یک مثال مناسب این است که به جای نگاه هم‌زمان به کل محیط، ابتدا با یک بخش تعامل شود و سپس توجه به بخش‌های دیگر منتقل گردد. کل محیط حذف نشده است؛ فقط توجه به‌صورت مرحله‌ای و محلی اعمال می‌شود. در مدل نیز \eng{local attention} به معنی حذف متن نیست، بلکه به معنی محدود کردن میدان توجه در هر گام و نیاز به ترکیب بعدی است.

\subsubsection{ترکیب پنجره‌های محلی}
اگر در هر پنجرهٔ محلی عناصر مهم پیدا شوند، هنوز باید مشخص شود چگونه این عناصر با عناصر مهم پنجره‌های دیگر ترکیب می‌شوند. یک راه ساده این است که از بازنمایی‌های خلاصه، میانگین‌ها یا توکن‌های تجمیعی استفاده شود. راه دقیق‌تر این است که لایهٔ بالاتری ساخته شود که بازنمایی‌های محلی را دوباره با هم مقایسه کند. بدون چنین مرحله‌ای، مدل ممکن است فقط چند نگاه محلی جداگانه داشته باشد و بازنمایی منسجم کل متن را نسازد.

\subsection{\eng{Multi-Head Attention} و \eng{Mixture of Experts}}
\eng{Multi-head attention} و \eng{Mixture of Experts} را نباید مستقیماً معادل راهکار حافظهٔ کاری دانست، اما هر دو می‌توانند به استفادهٔ بهتر از ظرفیت مدل کمک کنند.

در \eng{multi-head attention}، مدل چند مجموعهٔ جدا از \(Q\)، \(K\) و \(V\) دارد. هر \eng{head} می‌تواند جنبه‌ای متفاوت از ورودی را برجسته کند: یک سر ممکن است رابطه‌های نحوی را دنبال کند، سر دیگر رابطه‌های معنایی، سر دیگر ارجاع‌های موجودیت‌ها، و سر دیگر نشانه‌های موقعیت یا ساختار جمله را. سپس خروجی سرها با هم ترکیب می‌شود. این کار شبیه نگاه کردن از چند زاویه است، نه الزاماً بزرگ‌تر کردن حافظهٔ کاری.

\eng{Mixture of Experts} نیز معمولاً به این معناست که بخش‌هایی از مدل یا «متخصص‌ها» برای نمونه‌ها یا توکن‌های مختلف فعال می‌شوند. هدف اصلی آن افزایش ظرفیت مؤثر یا تخصصی کردن پردازش است، بدون اینکه همیشه همهٔ پارامترها فعال شوند. این ایده می‌تواند غیرمستقیم به مدیریت اطلاعات کمک کند، اما مسئلهٔ اصلی حافظهٔ کاری، یعنی نگه‌داری و بازیابی هدفمند اطلاعات زمینهٔ فعلی، همچنان به طراحی توجه، حافظه، خلاصه‌سازی و راهبردهای استدلال وابسته است.

\subsection{جمع‌بندی}
در این فصل دو مسیر برای کاهش محدودیت حافظهٔ کاری مدل‌ها از هم جدا شد. مسیر اول شناختی و الگوریتمی بود: \eng{chunking}، زیرهدف، خلاصه‌سازی و تقسیم مسئله. مسیر دوم معماری بود: بررسی اینکه در \eng{Transformer}، توجه چگونه بازنمایی زمینه‌مند می‌سازد و چه تغییرهایی مانند \eng{local attention} می‌توانند فشار محاسباتی و حافظه‌ای را کم کنند.

نکتهٔ مشترک هر دو مسیر این است که حافظهٔ مؤثر فقط به مقدار فضای خام بستگی ندارد. آنچه اهمیت دارد، سازمان‌دهی، انتخاب، بازیابی و فشرده‌سازی اطلاعات است.
